\documentclass[11pt,a4paper]{article}

\usepackage[margin=0.78in]{geometry}
\usepackage{setspace}
\usepackage{enumitem}
\usepackage{fancyhdr}
\usepackage{xcolor}
\usepackage{titlesec}
\usepackage{tcolorbox}

\definecolor{darkblue}{RGB}{35,65,105}

\setstretch{1.05}

\pagestyle{fancy}
\fancyhf{}
\lhead{\textbf{Yashjot Kaur}}
\rhead{\textcolor{darkblue}{Week 3 Journal}}
\cfoot{\thepage}

\renewcommand{\headrulewidth}{0.5pt}

\titleformat{\section}
{\large\bfseries\color{darkblue}}
{}{0em}{}

\begin{document}

\begin{center}
    {\LARGE \textbf{\textcolor{darkblue}{WEEK 3 JOURNAL}}}\\[5pt]
    {\large \textbf{AI-based Real-Time Sensitive Data Detection and Masking}}\\[6pt]
    \textbf{Name: Yashjot Kaur \quad | \quad Roll No.: 1024170452}
\end{center}

\vspace{4pt}
\hrule
\vspace{8pt}

\section*{1. Objective}

The main objective of Week 3 was to move from understanding the project
concept to starting the actual project work. I focused on understanding
how sensitive information can be detected from text and how the detected
information can later be masked.

\section*{2. Work Done}

During this week, I worked on the initial development of the project. I
created sample text containing information such as names, email addresses,
phone numbers and other personal details.

I studied how the system can identify such information using simple
patterns and NLP techniques. I also looked at how the detected information
can be separated into categories such as \textbf{EMAIL}, \textbf{PHONE},
and \textbf{PERSON}.

\begin{center}
\begin{tcolorbox}[
    colback=gray!5,
    colframe=darkblue,
    width=0.9\textwidth,
    arc=2mm,
    boxrule=0.5pt
]
\centering
\textbf{Text Input}
$\rightarrow$
\textbf{Find Sensitive Information}
$\rightarrow$
\textbf{Identify its Type}
$\rightarrow$
\textbf{Mask the Information}
\end{tcolorbox}
\end{center}

\section*{3. Practical Understanding}

I worked with different examples to understand how the same type of
information can appear in different formats. For example, phone numbers
may contain spaces, country codes or different separators.

This helped me understand why the detection rules need to be tested with
different inputs instead of checking only one fixed format.

\section*{4. Report Work}

Along with the technical work, I contributed to the preparation of the
project report. I worked on organising the content and documenting the
project background, objectives, methodology and the work completed so far.

I also helped in presenting the technical workflow in a clear and
structured manner so that the report explains the project development
step by step.

\section*{5. My Contribution}

My main contribution this week was understanding and working on the
initial detection stage. I also contributed to testing sample inputs and
documenting the progress of the project in the report.

This helped me connect the theoretical concepts studied earlier with the
actual project implementation.

\section*{6. Learning Outcomes}

This week helped me understand how a real project moves from an idea to
implementation. I learned that detecting sensitive information is not
only about finding patterns, but also about checking different cases and
improving the detection process.




\vfill

\begin{center}
    \textcolor{darkblue}{\rule{0.35\textwidth}{0.5pt}}\\[3pt]
   
\end{center}

\end{document}